\subsection{The marked arithmetic CM comparison and the actual congruence base}
\label{sec:arithmetic-cm-base}

The following results remove the arithmetic base hypothesis left explicit
in E.91--E.94.  They concern the actual Galois subgroup and ordinary
peripheral lifts.  They do not assume a sewing-compatible action of the
whole GT group in genus one.

Use \(\mathfrak B=C_2\amalg C_3=\langle s,r\rangle\),
\(\tau=sr\), \(F=\ker(\mathfrak B\to S_3)\), and
\(x=\tau^{-2}\), \(y=sxs\).  Put
\[
 U=\begin{pmatrix}1&1\\0&1\end{pmatrix},\quad
 L=\begin{pmatrix}1&0\\-1&1\end{pmatrix},\quad
 S=ULU,\quad R_0=-UL.
\]
Thus \(S^2=-I\), \(R_0^3=I\), and \(SR_0=U^{-1}\).
The signed pure representation sends \(x,y\) to \(U^2,L^2\);
the modular representation sends \(s,r\) to the classes of \(S,R_0\).
Write \(r_*=(3/4)^{1/4}>0\).

\begin{result}{E.95}{Proved}{The arithmetic comparison retains the two CM frames and their signs}
Let \(\sigma\in G_{\mathbb Q}\) have cyclotomic character one and
\(\kappa_2(\sigma)\equiv0\pmod3\), where \(\kappa_2\) is the
Kummer cocycle of the positive roots of two.  Let
\[
 \Phi_\sigma(s)=\beta_\sigma s\beta_\sigma^{-1},\qquad
 \Phi_\sigma(r)=A_\sigma rA_\sigma^{-1},
 \qquad A_\sigma,\beta_\sigma\in F.
\]
At modulus twelve the congruence intertwiner is the identity.
The signed matrices \(q_{12}(\beta_\sigma)^{-1}\) and
\(q_{12}(A_\sigma)^{-1}\), in the transported bases, are precisely
the Galois actions on the Legendre fibres at \(1/2\) and at
\(\lambda_0=(1+\sqrt{-3})/2\).  Both are identity modulo four, and
\[
 \epsilon_2\bigl(q_{12}(\beta_\sigma)^{-1}\bigr)
 =\epsilon_3\bigl(q_{12}(A_\sigma)^{-1}\bigr)
 =\begin{cases}0&\sigma(r_*)=r_*,\\1&\sigma(r_*)=-r_*.\end{cases}
                                                        \tag{E.95.1}
\]
Here \(\epsilon_2\) is exponent parity in
\(\langle S\bmod3\rangle=C_4\), and \(\epsilon_3\) is the sign on
\(\ker(R_0-I)=\mathbb F_3(1,-1)\).
In particular this identifies the exact arithmetic CM bits used
in LM4 A.31.  The comparison is not inferred from equality of
two division fields.

The historical finite discrepancy character \(\delta\) on the
multiplier-one gentle group is a continuous character, invariant
under conjugation by the full gentle group, and is zero on its
arithmetic subgroup.  Its extension \(\Delta\), obtained by
declaring the Galois image to have value zero, is therefore
well-defined.  This establishes the arithmetic input of recovered
LM7 E.65 without establishing its restriction to the full GT group.
\end{result}
\begin{proof}
\emph{The actual paths.}
Consider
\[
 \mathcal E_\lambda:\quad Y^2=X(X-1)(X-\lambda)
 \quad\hbox{over }\mathbb P^1-\{0,1,\infty\},
\]
based at the integral tangential point \(\overrightarrow{01}\).
The signed geometric monodromy, in the indicated convention, is
\((U^2,L^2)\).  The quotient orbifold by \(S_3\) has geometric
fundamental group \(\mathfrak B\).  Its order-two and order-three
stabilizers occur at \(1/2\) and \(\lambda_0\), respectively.
The primary comparison is explicit in
\citet[Proposition~2.19, p.437; \S4.14, equation~(4.15), p.445;
Theorem~4.16 and its proof, p.446; \S5, p.447]{SymNTY2010}.
In particular the latter discussion identifies the standard
Galois action on the quotient orbifold, rather than an arbitrarily
chosen action on an abstract modular group.

Choose a geometric path in the level-two base from the tangent to
each of these points.  It can be adjusted by a pure loop so that
the transported stabilizer is exactly \(s\), or exactly \(r\).
Indeed the torsion theorem for the full profinite free product
conjugates any such stabilizer to the corresponding free factor.
The image of the conjugator in \(S_3\) centralizes the named
transposition or three-cycle.  That centralizer is precisely its
cyclic subgroup; multiplying by a suitable factor element makes
the conjugator pure.  This uses only the free-factor case of
\citet[Theorem~B, pp.159--160]{SymHerfortRibes1985}.

The chosen \(\sigma\) fixes both level-two points and their
stabilizers, since \(\mathbb Q(\sqrt{-3},i)\subset\mathbb Q^{\rm ab}\).
If \(\gamma\) is the path, its Galois path cocycle is therefore
pure.  Transport of the stabilizer shows that its inverse is the
corresponding frame \(A_\sigma\) or \(\beta_\sigma\).
There is no further freedom: the centralizer of \(s\) or \(r\)
in \(\mathfrak B\) is that finite factor, whose intersection
with \(F\) is trivial.  This argument identifies the actual
frames, not just their conjugacy classes.

For a local system and a path \(\gamma:b\to p\), transport of
the arithmetic action gives
\[
  \rho_{p,\gamma}(\sigma)
       =\rho_{\rm geom}(\gamma^{-1}\sigma\gamma)\rho_b(\sigma).
                                                        \tag{E.95.2}
\]
The preceding frame convention makes the first factor the inverse
frame matrix.  This is an identity of fibre functors; at finite
level it follows by applying \(\sigma\) to the path isomorphism
on the finite torsion cover.  Hence it retains the signed matrix.

\emph{Why the tangential factor is one here.}
The tangential Tate-module matrix is
\(\left(\begin{smallmatrix}\chi&8\kappa_2\\0&1\end{smallmatrix}\right)\)
in the standard Ferrari convention
\cite[p.446, proof of Theorem~4.16]{SymNTY2010}.
Reversing a homology convention changes the shear sign, which
does not affect its vanishing in the present subgroup.
The Legendre calculation independently gives the same vanishing:
its \(j\)-invariant is
\[
  \frac{256(1-\lambda+\lambda^2)^3}
       {\lambda^2(1-\lambda)^2},
\]
so its formal Tate parameter is
\(q=\lambda^2\,2^{-8}(1+\lambda\mathbb Q[[\lambda]])\).
Every finite torsion level is generated by the roots of unity,
a root of this \(q\), and the splitting quadratic twist.
The unit has its unique root with constant coefficient one;
the splitting twist has constant square class \(-1\).
Thus for \(\chi=1\) the only possibly nontrivial coefficient
is a sign convention times \(8\kappa_2\); the twist is already
fixed.  For the Puiseux splitting and the actual Tate curve
comparison see \citet[\S5.2, pp.449--450; \S5.7, pp.456--458;
Proposition~5.7.3 stated on p.458 and proved on pp.481--482]
{SymNakamura2013}.
Our hypothesis makes this shear zero modulo twelve.
Equation (E.95.2) proves the asserted fibre comparison.

\emph{The harmonic fibre.}
The substitution \(X'=4X-2,\ Y'=8Y\) gives
\[
             \mathcal E_{1/2}:\quad (Y')^2=(X')^3-4X',
\]
the quadratic twist by two of \(E_4:Y^2=X^3-X\).
The twist is trivial on \(G_{\mathbb Q^{\rm ab}}\).
The four-division points are defined over
\(\mathbb Q(i,\sqrt2)\): halves of \((0,0)\) have coordinates
\((\pm i,\pm(1\mp i))\), and halves of \((1,0)\) have
\((1\pm\sqrt2,\pm(2\pm\sqrt2))\).
Their doubles generate \(E_4[2]\), so these halves generate
\(E_4[4]\).  The harmonic matrix is consequently \(I\) modulo four.

Put
\[
 a=\frac{r_*^{-1}+2r_*}{2},\qquad
 b=\frac{r_*^{-1}-2r_*}{2i}.
\]
The roots of the third-division polynomial \(3X^4-6X^2-1\)
are \(\pm a,\pm b\).  Since \(\sigma\) fixes \(\sqrt3=2r_*^2\)
and \(i\), it sends \(r_*\) to \(r_*\) or \(-r_*\).
In the first case it fixes each third-division \(X\)-coordinate.
In the second it sends every such coordinate to its negative,
which is the projective action of the order-four CM automorphism.
The action on \(E_4[3]\) commutes with that automorphism; its
centralizer in \(\operatorname{SL}_2(\mathbb F_3)\) is \(C_4\).
The kernel of its action on the four pairs \(\{P,-P\}\) is
\(\{\pm I\}\).  These facts prove the first bit in (E.95.1).
The transported order-four automorphism is \(S\) or \(S^{-1}\);
its exponent parity is the same in either orientation.

\emph{The equianharmonic fibre and its necessary twist.}
Writing \(X=Z+(1+\lambda_0)/3\) gives
\[
 \mathcal E_{\lambda_0}:\quad Y^2=Z^3+d^3,\qquad
 d=-\frac{2\lambda_0-1}{3}=\frac1{\sqrt{-3}}.
                                                        \tag{E.95.3}
\]
It is the quadratic twist by \(d\) of \(E_3:Y^2=X^3+1\).
On \(G_{\mathbb Q^{\rm ab}}\) the twist character is
\[
 \frac{\sigma(\sqrt d)}{\sqrt d}
   =\frac{\sigma(r_*)}{r_*},\qquad
 (-3)^{1/4}=\zeta_8\sqrt2\,r_*.
                                                        \tag{E.95.4}
\]
The inverse arising from \(\sqrt d=(-3)^{-1/4}\) is immaterial
because this character is a sign.
The point
\[
 P=(-1+\sqrt3,\ \sqrt3(\sqrt3-1)r_*)
\]
has double \((-1,0)\); \(P\) and its \(\zeta_3\)-multiple
form a basis of \(E_3[4]\).  Therefore \(\sigma\) acts on
this entire four-torsion module by
\((\sigma(r_*)/r_*)I\).  The twist in (E.95.4) cancels that
sign, giving \(I\) modulo four on the Legendre fibre.
On the other hand \((0,1)\in E_3[3](\mathbb Q)\) spans the
fixed line of its order-three CM automorphism.  Its action
after twisting is exactly the sign in (E.95.4).  The transported
line is \(\ker(R_0-I)\); replacing the CM generator by its inverse
does not change this line.  This proves the second bit.
Omitting this twist would give the wrong signed statement.

\emph{The finite character and its arithmetic restriction.}
For completeness, this portion of recovered LM7 E.65 can be
reconstructed without a four-strand sewing assertion.
For a positive gentle element \(g\), write
\((\beta_g^{-1})^{\rm ab}=(\rho_g,\rho_g)\).
The class-two hexagon gives zero central class of \(f_g\);
the identity \(f_g=\theta(\beta_g^{-1})^{-1}\beta_g^{-1}\),
where \(\theta\) interchanges \(x,y\), then gives
\(\rho_g^2+2e=0\).  Thus \(\rho_g\in2\widehat{\mathbb Z}\).
Modulo four the signed pure image is the four-group generated
by \(U^2,L^2\), and does not contain \(-I\).
It follows that \(q_4(\beta_g)=I\).

At three put \(X_0=U^2\).  Direct multiplication in the
twenty-four-element group \(\operatorname{SL}_2(\mathbb F_3)\)
gives
\[
 \bigl((Sg_0S^{-1}g_0^{-1})^{-1}R_0\bigr)^3=I
 \ \Longleftrightarrow\
 g_0=X_0^kS^j\quad(0\le k<3,\ 0\le j<4).
                                                        \tag{E.95.5}
\]
The other twelve matrices have cube \(-I\).
The quotient character to \(C_3\) sends both \(U^2,L^2\) to
one, so \(k=-2\rho_g=\rho_g\pmod3\).
Combining this calculation with the reduction at four gives
\[
 q_{12}(f_g)=
 \begin{pmatrix}1&-8\rho_g\\-8\rho_g&1+64\rho_g^2\end{pmatrix},
 \qquad V_g=U^{8\rho_g},\qquad q\Phi_g=\operatorname{Ad}(V_g)q.
                                                        \tag{E.95.6}
\]
The \(R_0\)-centralizer in the signed pure four-group is trivial,
so \(q_4(A_g)=I\) as well.  Consequently
\[
 T_g=q_{12}(A_g)^{-1}V_g,\qquad
 W_g=q_{12}(\beta_g)^{-1}V_g
                                                        \tag{E.95.7}
\]
are \(I\) at four and belong at three to the Cartans \(C_6\)
and \(C_4\), respectively.
The pure-frame law
\(A_{gh}=\Phi_g(A_h)A_g\), and the identical law for \(\beta\),
together with (E.95.6), proves that \(T,W\) are homomorphisms.
Thus \(\delta=\epsilon_3(T)+\epsilon_2(W)\) is a continuous
character.

For arbitrary multipliers, the braid relation and trace two
give the unique congruence intertwiner
\(\left(\begin{smallmatrix}\lambda&t\\0&1\end{smallmatrix}\right)\)
prime by prime.  Indeed solving the relation with
\(U^\lambda\) gives either the abelian pair, excluded by
surjectivity, or
\(\left(\begin{smallmatrix}1-u&\lambda u^2\\-1/\lambda&1+u
\end{smallmatrix}\right)\).
This is the elementary calculation of
\citet[Lemmas~4.1--4.3, pp.774--776]{SymMarin2008};
it uses only the actual three-strand action.
Its two transported frame matrices normalize the corresponding
Cartans.  At three, a normalizer preserves the fixed line of
the \(C_3\)-part, and sends the order-four generator to itself
or its inverse.  Thus both bits, and \(\delta\), are invariant
under these conjugations.  The sign is fixed by reduction at four.

On the arithmetic multiplier-one subgroup, the kernel of
\(\kappa_2\bmod3\) has index dividing three.
Its actual congruence shear is zero modulo twelve; uniqueness of
the normalized shear identifies it with \(V_g\) in (E.95.6).
The marked comparison already proved shows \(\delta=0\) on this
kernel.  A homomorphism from a group of order dividing three to
\(\mathbb F_2\) is zero, so \(\delta\) vanishes on the whole
arithmetic multiplier-one subgroup.
Cyclotomic surjectivity now defines
\(\Delta(g)=\delta(g\sigma^{-1})\) for any arithmetic
\(\sigma\) with \(\chi(\sigma)=\lambda(g)\).
The proved vanishing makes this independent of the choice;
conjugation invariance proves multiplicativity.
The compact quotient map from the positive group times
\(G_{\mathbb Q}\) proves continuity.
\end{proof}

\begin{result}{E.96}{Proved}{An arithmetic involution with zero complete Kummer shear and both nonzero CM bits}
There is \(\sigma\in G_{\mathbb Q}\) satisfying
\[
 \chi(\sigma)=1,\quad \kappa_2(\sigma)=0,\quad
 c\sigma c=\sigma^{-1},\quad \sigma(r_*)=-r_*,
                                                        \tag{E.96.1}
\]
such that \(\sigma c\) is a Galois conjugate of complex conjugation.
Its action is the identity on every congruence quotient \(Q_N\).
For its signed pure frames at twelve,
\[
 q_{12}(A_\sigma)=5I,\qquad
 q_{12}(\beta_\sigma)=B^{\pm1},\qquad
 B=\begin{pmatrix}9&4\\8&9\end{pmatrix}.
                                                        \tag{E.96.2}
\]
\end{result}
\begin{proof}
Fix an embedding in \(\overline{\mathbb Q}_3\).
Over the unramified quadratic field \(\mathbb Q_3(i)\),
\(4X^4-3\) is Eisenstein and adjoining \(r_*\) is a cyclic
tamely ramified extension of degree four.  Choose an absolute
inertia lift \(\gamma\) with \(\gamma(r_*)=ir_*\).
All splitting fields of \(X^{2^n}-2\) and all dyadic cyclotomic
fields are unramified at three, as their discriminants have
only the prime two.  Thus \(\gamma\) fixes the positive
dyadic roots of two and all dyadic roots of unity.
Let \(\sigma_0=\gamma c\gamma^{-1}c\).
Its cyclotomic character is one, its dyadic Kummer character
of two is zero, \(c\sigma_0c=\sigma_0^{-1}\), and direct
application to \(r_*\) gives \(-r_*\).

Let \(e_2\in\widehat{\mathbb Z}\) have dyadic component one
and all odd components zero; take
\(\sigma=\sigma_0^{e_2}\) in its procyclic closure.
The Kummer cocycle is additive on \(\ker\chi\).
This kills every odd component and retains the zero dyadic
component, proving (E.96.1).
For \(b=(e_2-1)/2\in\widehat{\mathbb Z}\), the identity
\[
 \sigma c=(\sigma_0^b\gamma)c(\sigma_0^b\gamma)^{-1}
                                                        \tag{E.96.3}
\]
proves the stronger assertion about complex conjugation.
The actual congruence shear is zero
\cite[\S4, p.776]{SymMarin2008}; hence every \(q_N\Phi_\sigma=q_N\).

Apply E.95 directly to \(\sigma\), which still moves \(r_*\).
The frame maps on the subgroup with zero congruence action are
homomorphisms into their Cartans.  On the procyclic image of
\(\sigma\), all finite quotients are two-groups.
Its \(C_6\)-value is therefore the unique nontrivial element
of order two, while its \(C_4\)-value is either odd generator.
Their unique signed lifts equal to \(I\) modulo four are
precisely (E.96.2).
\end{proof}

\begin{result}{E.97}{Proved}{The historical positive towers have an actual marked congruence base}
For every \(12\mid N\) put
\[
 U_N=\ker(\mathfrak B\to Q_N),\qquad
 \Gamma_N=\mathfrak B/R_{\rm nil}(U_N).
\]
There is one ordinary full positive gentle element \(h\) such that
\(a=hc\) has parity \(\Delta(a)=1\), and its actual images on
all \(\Gamma_N\) are compatible involutions with literal cusp
inversion and both involutory peripheral normalizers.
Their outer images on \(F/R_{\rm nil}(U_N)\) also have order two,
and the parity character descends to these images.

Writing \(K_{\rm cong}=\bigcap_NU_N\), the resulting data on
\[
 C=\mathfrak B/R_{\rm nil}(K_{\rm cong})
                                                        \tag{E.97.1}
\]
can be chosen with the precise head projections required before
E.91: positive action identity on
\(H=\mathbb F_7^{13}\rtimes Q_6\), positive \(s\)-frame
\((0,\beta_0)\), positive \(r\)-frame one, and negative
\(r\)-conjugator \(s\).
Thus E.91--E.94 apply to actual data.  This confirms the
arithmetic base and compatible-frame assertions of LM4 A.31/A.33
and the positive constructions used in A.76/A.79/A.81.
It does not make the relative tower cofinal or make \(a^2=1\)
in the full group.
\end{result}
\begin{proof}
\emph{The relative components and the required elementary lifts.}
For a finite quotient \(D\) retaining the two factors, let
\(\Gamma_p=\mathfrak B/O^p(\ker q)\), with pro-\(p\) kernel.
If \(u=s\) or \(r\) has order prime to \(p\), then
\[
 C_{\Gamma_p}(u)\longrightarrow C_D(q(u))
                                                        \tag{E.97.2}
\]
is onto.  Lift a finite centralizer element; the two resulting
copies of \(\langle u\rangle\) are complements in the same
pro-\(p\) extension.  Coprime complement conjugacy adjusts
the lift to centralize \(u\), first in every finite quotient
and then by compactness.
At the matching prime, \(C_{\Gamma_p}(u)=\langle u\rangle\).
Indeed the preimage of \(\langle q(u)\rangle\) is an open
free profinite product of the corresponding \(C_p\)-factors
and a free factor.  Its relative pro-\(p\) quotient is its
maximal pro-\(p\) quotient.  The distinct factor labels remain
nonconjugate there, and each factor has its own centralizer
\cite[Theorem~B', pp.159--160]{SymHerfortRibes1985}.

We also need the following prescribed version of a splitting.
Every pro-two extension of \(S_3\) having an involution over a
named transposition has a section retaining that involution.
For any finite \(\mathbb F_2[S_3]\)-module \(M\), restriction
\[
 H^j(S_3,M)\longrightarrow H^j(C_2,M)\quad(j>0)
                                                        \tag{E.97.3}
\]
is an isomorphism.  Restriction after corestriction has one
identity double-coset term and one term through the trivial
intersection; corestriction after restriction multiplies by
the odd index three.  Both composites are identity.
For a finite two-kernel, induct on its order using the
characteristic subgroup of central elements killed by two.
Degree two gives a splitting because the prescribed reflection
splits the restriction; degree one and a kernel conjugation
make the lifted reflection the prescribed one.
Compactness of the finite tuples defining a section proves
the assertion for a pro-two kernel.  Thus no projectivity of
\(S_3\) has been assumed.

\emph{The actual action.}
Take the single arithmetic \(\sigma\) of E.96.  The fibre product
of the groups \(\Gamma_{N,p}\) over \(Q_N\) is \(\Gamma_N\).
Define \(\phi_N\) by \(\Phi_\sigma\) on its dyadic factor and
by identity on every odd factor.
The common finite action is identity, so these are actual
compatible automorphisms; \((\phi_Nc)^2=1\).
Let \(u_N=q_N(A_\sigma)\) and \(b_N=q_N(\beta_\sigma)\).

\emph{The dyadic \(r\)-frame with its exact reflection.}
Use signed matrices \(R_0\) above and
\(D=\left(\begin{smallmatrix}0&1\\1&0\end{smallmatrix}\right)\).
The matrix \(u_N\) centralizes \(R_0\) and is \(I\) at four.
For any determinant-one element \(u=vI+wR_0\) of this Cartan,
\[
 DuD=u^{-1},\qquad (u^{-1}D)^2=I,\qquad
                  \operatorname{tr}(u^{-1}D)=0.
                                                        \tag{E.97.4}
\]
Over \(\mathbb Z_p\), \(u^{-1}D\) is conjugate to \(D\) by
a special-linear matrix, pure modulo two.
For odd \(p\), its two eigensummands are rank-one direct
summands; adjust the determinant of an eigenbasis using the
centralizer of \(D\).
At two put \(T=u^{-1}D\) and \(H_0=(e_1,Te_1)\).
Then \(H_0DH_0^{-1}=T\), \(H_0\equiv I\pmod2\), and
\(d_0=\det H_0\equiv1\pmod4\).
Multiplication on the right by
\[
 \begin{pmatrix}v&w\\w&v\end{pmatrix},\qquad
 v=(d_0^{-1}+1)/2,\quad w=(d_0^{-1}-1)/2
                                                        \tag{E.97.5}
\]
corrects the determinant to one while retaining both properties.
Reduction and the Chinese remainder theorem give a pure
conjugator at every \(N\); it lifts to \(F\).

In the dyadic extension by the genuine involution
\(a_2=\Phi_\sigma c\), set
\[
 J_2=\Phi_\sigma(s)A_{\sigma,2},\qquad
 D_2=J_2^{-1}a_2,\qquad T_2=\beta_{\sigma,2}^{-1}a_2.
                                                        \tag{E.97.6}
\]
Already in the full group \(D_2^2=1\): its square lies in
\(\langle r\rangle\) and is fixed by the inversion which
\(D_2\) induces there.  Also \(T_2^2=1\): its square lies
in \(\langle s\rangle\), but cannot be \(s\), whose image is
nonzero in the exponent-two abelianization of the extension.
The finite image of \(D_2\) is \(u_N^{-1}d\), \(d=sc\).
The matrix calculation conjugates it purely to a reflection
with finite image \(d\).  Apply (E.97.3) in the preimage of
\(\langle q_N(r),d\rangle\simeq S_3\), retaining that reflection.
Coprime conjugacy then makes its three-cycle the named \(r\),
without changing the finite image.  Denote its reflection by
\(D'_2\).  The element \(D_2D'_2\) centralizes \(r\), has
finite image \(u_N^{-1}\), and changes the positive \(r\)-frame
to one of finite image one with involutory negative normalizer.
At every odd prime take positive \(r\)-frame one and negative
normalizer \(sc\).

\emph{The \(s\)-frame and the specified seven-primary head.}
The signed \(b_N\) centralizes \(S\).
Conjugation by \(\operatorname{diag}(-1,1)\) inverts its
norm-one Cartan.  Thus \(b_N^{-1}c\) is a finite involution.
At two retain the genuine \(T_2\).  At any odd prime,
(E.97.2) and coprime splitting lift this finite involution
inside \(C_{\Gamma_{N,p}}(s)\rtimes\langle c\rangle\).
Its associated positive \(s\)-frame has finite image \(b_N\).
All frames are pure because those finite images are pure.

The extra head prescription is compatible with this argument.
The actual head \(H=\mathfrak B/\overline{K_6'K_6^7}\)
has \(K_6=\ker(\mathfrak B\to Q_6)\); its seven-kernel is
\(N_7=\mathbb F_7^{13}\).  Every solvable quotient of a
congruence quotient has no \(C_7\) composition factor, so the common quotient of \(H\)
and \(Q_N\) is exactly \(Q_6\).  Their actual joint quotient is
therefore \(H\times_{Q_6}Q_N\), and it is a quotient of
\(\Gamma_{N,7}\).
In the source coordinates of E.59, \(\beta_0\) centralizes
the named \(s\), has order two, and is fixed by \(c\).
The pair \((b_N^{-1}c,\beta_0c)\) is consequently an actual
involution in that joint quotient and centralizes \(s\).
Applying the same centralizer and coprime lifting argument
over this joint quotient, instead of just \(Q_N\), gives the
literal head \((0,\beta_0)\).
The positive action is already identity on \(H\), and the
retained odd \(r\)-frame is one.  Thus its negative
\(r\)-conjugator is the actual source \(s\), as required.

\emph{One ordinary full lift and the inverse limit.}
The relative kernel is relatively perfect, so E.57 realizes
each specified action and its literal frames by a full
peripheral lift.  Normalize it to fix \(\tau\) literally.
Its commutator normalization is retained: the quotient
\(f\)-coordinate has zero abelianization at every prime, and
\(R_{\rm nil}(U_N)\subset F'\).  Thus it is a positive
gentle lift, not merely an unnormalized automorphism.
The two involutory-normalizer equations say that the
frames of its negative square are identity in \(\Gamma_N\).
All these are closed conditions in the compact space of
normalized full lifts and literal conjugators.
A common multiple of any finite set of moduli supplies
all its requirements, including the fixed head.
Compactness gives one \(h\) satisfying them for every \(N\).
No surjectivity between chosen frame-solution sets is used.

The parity at twelve is computed without a new arithmetic
comparison.  Its \(r\)-frame has image one, its \(s\)-frame
has image \(B^{\pm1}\), and its congruence action is identity.
Thus \(V_h=I\), and (E.95.7) gives \(\delta(h)=1\).
Since \(\Delta(c)=0\), \(\Delta(hc)=1\).
Its multiplier is \(-1\), which acts as \(-I\) on the full
abelianization \(\widehat{\mathbb Z}^2\); hence each outer
image has order exactly two.

Finally
\[
 R_{\rm nil}(K_{\rm cong})
       =\bigcap_{12\mid N}R_{\rm nil}(U_N).              \tag{E.97.7}
\]
To prove this, pass to any finite source quotient.
The descending images of \(U_N\) stabilize at the image
of their intersection, by compactness.  Their nilpotent
residuals stabilize at that of this finite image.
The image of the nilpotent residual of any profinite
subgroup is the residual of its finite image: both are
the stable intersection of the lower central terms.
Finite separation proves (E.97.7).
The data just constructed therefore live on the exact
group \(C\) of (E.97.1).

\emph{Descent to the outer image.}
At modulus twelve the discrete modular cover has forty-eight
cusps.  In its pro-\(p\) first homology their oriented vectors
have just their sum as relation, so the stabilizer of the
\(\tau^{12}\)-vector is precisely \(\langle q_{12}(\tau)\rangle\).
Since this boundary element is primitive in the free kernel,
the free pro-\(p\) centralizer theorem implies
\(C_{\Gamma_{12}}(\tau)=\overline{\langle\tau\rangle}\);
the same argument applies to \(x=\tau^{-2}\), and conjugation
by \(s\) treats \(y\).
Their centralizers intersect trivially: their finite \(S_3\)
images are distinct order-two groups, and in the pure
subgroup the two abelian axes have zero intersection.
Thus \(C_{\Gamma_{12}}(F_{12})=1\).

If a normalized action is outer-trivial on \(F_{12}\),
its multiplier is one and it is \(\operatorname{Inn}(u)\)
on \(F_{12}\) for \(u\in F_{12}\).
The preceding centralizer makes it the same inner action
on \(\Gamma_{12}\).  Literal cusp fixation gives \(u=x^a\).
Its formula at \(s\) then gives \(f=y^ax^{-a}\);
the normalized condition \(f\in F_{12}'\) forces \(a=0\).
An identity actual action has trivial finite \(s\)-frame
by the matching-prime centralizer at two and trivial
\(r\)-frame by that at three.  Its shear is also identity.
Hence its discrepancy is zero, proving descent of
\(\Delta\) to the outer image.
Every larger \(N\) maps to this outer quotient.
The noncofinality remains the independent \(A_7\)-quotient
obstruction E.56; nothing here asserts a full involution lift.
\end{proof}
